\documentclass[11pt]{article}

\usepackage[T1]{fontenc}
\usepackage[utf8]{inputenc}
\usepackage{libertine}
\usepackage[libertine]{newtxmath}
\usepackage[margin=1.1in]{geometry}
\usepackage{microtype}
\usepackage{amsmath}
\usepackage{booktabs}
\usepackage{graphicx}
\usepackage[numbers,sort&compress]{natbib}
\usepackage[dvipsnames]{xcolor}
\usepackage[colorlinks=true,linkcolor=MidnightBlue,citecolor=MidnightBlue,urlcolor=MidnightBlue]{hyperref}
\usepackage[capitalise,noabbrev]{cleveref}
\usepackage{enumitem}
\setlist{nosep,leftmargin=*}

\newcommand{\pending}{\textcolor{BrickRed}{[pending]}}
\newcommand{\dnll}{\Delta\mathrm{NLL}}
\newcommand{\op}[1]{\textsf{#1}}

% Numbers come from probe.py -> oracle_out_full/results.json -> scripts/make_tables.py.
% Until that has run, every result renders as [pending].
\IfFileExists{generated/results.tex}{\input{generated/results.tex}}{% Fallback used until scripts/make_tables.py writes generated/results.tex.
\newcommand{\resPositions}{\pending}
\newcommand{\resBaseNLL}{\pending}
\newcommand{\resOracleGain}{\pending}
\newcommand{\resOraclePct}{\pending}
\newcommand{\resFracHelped}{\pending}
\newcommand{\resBestStatic}{\pending}
\newcommand{\resBestStaticDnll}{\pending}
\newcommand{\resSafeSkip}{\pending}
\newcommand{\resLagRecovery}{\pending}
\newcommand{\HeadroomRows}{\op{repeat} / \op{skip} / \op{swap} / \op{block} & \pending & \pending & \pending \\}
\newcommand{\RouterRows}{\pending & & & & \\}
\newcommand{\AblationRows}{\pending & & \\}
\newcommand{\BudgetTable}{\pending}
\newcommand{\SAENote}{\pending}
}

\title{\textbf{Layer Looping}\\[4pt]
\large Per-token layer programs in a frozen LLM:\\ how much is there to gain, and can the previous token see it coming?}
\author{Tyler Gee\\ \small K\=aula}
\date{\small Preprint draft \textperiodcentered{} \today}

\begin{document}
\maketitle

\begin{abstract}
\noindent
A pretrained transformer sends every token through every layer exactly once. We ask how much
next-token loss is left on the table by that choice, and whether the right per-token change can
be chosen \emph{before} the token is processed. For every token of a frozen Gemma-2-2B we compute
the exact change in its next-token negative log-likelihood under 104 layer programs: repeat one
layer $k$ times, skip one layer, swap two adjacent layers, or loop the block of layers 10--13. Only
that token's computation changes. From this per-token oracle we measure (i)~headroom by edit
type and layer; (ii)~a signed-cost compute--accuracy frontier on which skips can pay for loops;
(iii)~how much of that headroom a linear router recovers from the \emph{previous} token's
internals, against a same-token router and training-free entropy- and norm-gated baselines from
prior work; and (iv)~whether interpretable signals (logit-lens entropy and agreement, residual
update norms, attention entropy, Gemma Scope SAE latents) carry that information as well as raw
hidden states. We also measure how much of the per-token gain survives when routes are applied to
whole sequences, where edited tokens change what later tokens attend to.
\end{abstract}

\section{Introduction}

Adaptive-depth methods---adaptive computation time~\citep{graves2016act}, universal
transformers~\citep{dehghani2019universal}, early exiting~\citep{schuster2022calm},
mixture-of-depths~\citep{raposo2024mod}, layer dropping~\citep{fan2020layerdrop,men2024shortgpt}---assume
tokens differ in how much computation they need. Frozen models tolerate having layers skipped,
repeated, or reordered~\citep{sun2024painters,lad2024robustness}, and recent training-free methods
re-apply layers per token at high-entropy positions~\citep{bhattacharjee2026eglr} or where the
hidden-state norm peaks~\citep{zhang2026telltale}. Static mid-block loops in frozen models give
gains within noise~\citep{lys2026innerloop,chen2026trainingfreeloop}, and learned routers have
repeatedly recovered little of oracle headroom~\citep{glavas2024dynamic,alizadeh2024duo}.

What is missing is a clean measurement: for each token in a frozen model, \emph{exactly} how much
would a small edit to its own layer schedule help, and is that predictable from information
available before the token's forward pass? Our contributions:
\begin{enumerate}
  \item \textbf{An exact per-token oracle} over repeat, skip, swap and block-loop programs in
        Gemma-2-2B, computed for all positions in about one forward pass per program (\cref{sec:oracle}).
  \item \textbf{A signed-cost frontier} in which skips ($-1$) pay for repeats, against static
        programs and training-free baselines from prior work.
  \item \textbf{Routing from token $t-1$}, decided before token $t$'s pass, compared with a
        same-token router.
  \item \textbf{A controlled comparison} of interpretable signals, raw hidden states and SAE latents
        as router inputs.
  \item \textbf{The interaction gap}: how much of the summed per-token gain survives when whole
        routes are applied to sequences.
\end{enumerate}

\section{Method}
\label{sec:oracle}

\paragraph{Layer programs.}
Gemma-2-2B has $L=26$ layers, applied in order $0,\dots,L-1$. A program $c$ is one edit to this
schedule for a single token: \op{repeat}@$\ell$\,$\times k$ ($k\in\{2,3\}$, cost $k-1$),
\op{skip}@$\ell$ (cost $-1$), \op{swap}@$\ell$ (run $\ell+1$ before $\ell$, cost $0$), plus the
static block loop of layers 10--13 run twice, the best block reported by \citet{lys2026innerloop}
(cost $4$). That is 104 programs. Cost counts extra layer applications.

\paragraph{The diagonal swap.}
We want token $t$'s next-token loss when only $t$ runs program $c$ and every other token runs
normally. Because the model is causal, we can compute this for all positions at once. At each
step of the program, a query attends to the \emph{base-pass} keys and values of every other
position at the layer being run, and to its \emph{own} modified key and value on the diagonal.
This is exact for each token's own loss, at the cost of about one forward pass per program. It
ignores the effect of $t$'s modified key and value on later tokens; \cref{sec:interaction}
measures that separately. CPU tests on tiny random Gemma-2 models check that every oracle delta
equals an independent recompute that truncates the sequence at $t$ and runs the program normally.

\paragraph{Oracle and routers.}
$\dnll_t(c) = \mathrm{NLL}_t(c) - \mathrm{NLL}_t(\text{base})$; negative means the edit helps. The
oracle picks $\arg\min_c \dnll_t(c)$, falling back to the base program when nothing helps. A token
is \emph{helped} if the oracle gain exceeds $0.05$ nats, and a skip is \emph{safe} if its
$\dnll < 0.01$. A router predicts the vector $\widehat{\dnll}_t$ by ridge regression and routes to
$\arg\min_c [\widehat{\dnll}_t(c) + \lambda\,\mathrm{cost}(c)]$ (base if the minimum is
non-negative); sweeping $\lambda$ traces its frontier, and we always report \emph{realized} $\dnll$.
The \textbf{same-token} router reads token $t$'s base-pass features; the \textbf{lagged} router
reads token $t-1$'s, which exist before $t$ is processed.

\paragraph{Features and baselines.}
Interpretable features are four per-layer groups: logit-lens~\citep{nostalgebraist2020logitlens}
entropy, logit-lens agreement with the final prediction, relative residual-update norm, and
attention entropy. We compare them with raw hidden states and with Gemma Scope
SAE~\citep{lieberum2024gemmascope} latents at layers 6, 12 and 18. Training-free baselines: an
EGLR-like rule~\citep{bhattacharjee2026eglr} that repeats a top layer when next-token entropy is
high (decided at $t$, or one token early), and a norm-peak rule~\citep{zhang2026telltale} that
repeats the layer with the largest update norm.

\paragraph{Data and evaluation.}
WikiText-103 prose and GSM8K question--answer text in alternating 512-token chunks, each starting
with BOS. Chunks with index $\equiv 0 \pmod 5$ are held out; 95\% confidence intervals come from
bootstrap resampling over held-out chunks.

\section{Results}

\subsection{Headroom (E1)}
Over \resPositions{} positions with base NLL \resBaseNLL{} nats/token, the per-token oracle reduces
NLL by \resOracleGain{} nats (\resOraclePct\% of base) and helps \resFracHelped{} of tokens. The
best static program, applied to every token, is \resBestStatic{} with mean $\dnll$
\resBestStaticDnll. A safe skip exists for \resSafeSkip{} of tokens.

\begin{table}[h]
\centering
\caption{Oracle headroom by edit type. \emph{Restricted gain}: oracle gain using only this edit
type. \emph{Best-op share}: fraction of all tokens that are helped and whose best edit is this type.}
\label{tab:headroom}
\begin{tabular}{lccl}
\toprule
Edit & Restricted gain (nats) & Best-op share & Best static \\
\midrule
\HeadroomRows
\bottomrule
\end{tabular}
\end{table}

\subsection{Compute--accuracy frontier (E2)}

\begin{figure}[h]
\centering
\IfFileExists{figures/frontier.png}{%
  \includegraphics[width=0.85\linewidth]{figures/frontier.png}%
}{%
  \fbox{\parbox[c][2.2in][c]{0.8\linewidth}{\centering\pending\\[2pt]\small\texttt{bash run\_experiments.sh full}, then \texttt{scripts/make\_tables.py}}}%
}
\caption{NLL saved versus extra layer applications per token. Curves are Pareto fronts over
$\lambda$; markers are static programs.}
\label{fig:frontier}
\end{figure}

\begin{table}[h]
\centering
\caption{Budget-matched $\dnll$ (nats/token; lower is better): the best frontier point using at most
the given extra layer applications per token. ``---'': unreachable.}
\label{tab:budget}
\small
\BudgetTable
\end{table}

\subsection{Does the previous token know? (E3)}
The lagged interpretable router retains \resLagRecovery{} of the same-token router's gain.

\begin{table}[h]
\centering
\caption{Routers at $\lambda=0$ on held-out chunks. AUC: detecting helped tokens from the predicted
best gain. \emph{Edit\%}: share of tokens routed away from the base program.}
\label{tab:routers}
\small
\begin{tabular}{lcccc}
\toprule
Features & $\dnll$ & 95\% CI & AUC & Edit\% / +layers \\
\midrule
\RouterRows
\bottomrule
\end{tabular}
\end{table}

\subsection{What carries the signal? (E4)}
\SAENote

\begin{table}[h]
\centering
\caption{Lagged routers using one interpretable feature group ($\lambda=0$).}
\label{tab:ablation}
\begin{tabular}{lcc}
\toprule
Feature group & $\dnll$ & +layers \\
\midrule
\AblationRows
\bottomrule
\end{tabular}
\end{table}

\subsection{Interaction gap (E5)}
\label{sec:interaction}
\pending{} Whole-sequence routed loss versus the sum of independent per-token oracle deltas, from
\texttt{routed\_eval.py} with first-pass KV caching.

\section{Related work}
\textbf{Per-token computation in frozen models.} EGLR~\citep{bhattacharjee2026eglr} re-applies the
top layers at high-entropy tokens and the Tell-Tale Norm~\citep{zhang2026telltale} recurses layers
guided by hidden-state norms; both decide from the current token, after its pass. Think-at-Hard~\citep{fu2025thinkathard}
adds whole-model iterations for hard tokens with a trained backbone. \textbf{Programs of layers.}
CoLa~\citep{li2025cola} searches per-sample skip/repeat chains, POLAR~\citep{li2026polar} learns
programs of layers, and Dr.LLM~\citep{heakl2025drllm} trains sequence-level routers. \textbf{Static
loops.} Inner Loop Inference~\citep{lys2026innerloop} and Training-Free Looped
Transformers~\citep{chen2026trainingfreeloop} loop fixed blocks. \textbf{Routers versus oracles.}
\citet{glavas2024dynamic} find per-token skip controllers no better than a constant input, and
Duo-LLM~\citep{alizadeh2024duo} finds large gaps between oracles and learned routers.
\citet{wei2026earlyexit} argue skip headroom is shrinking in modern models. \textbf{Adaptive
depth with training.} ACT~\citep{graves2016act}, universal transformers~\citep{dehghani2019universal},
CALM~\citep{schuster2022calm}, mixture-of-depths~\citep{raposo2024mod}.

\section{Limitations}
The oracle is a hindsight upper bound: it picks the best of 104 outcomes per token, which rewards
any edit that perturbs a wrong prediction. Per-token deltas ignore effects on later tokens; E5
measures that gap. Routers are linear, the model is a single 2B checkpoint, and loss is measured
with teacher forcing.

\bibliographystyle{plainnat}
\bibliography{references}

\end{document}
